\section{Полугласные}

В румынском языке два полугласных звука: \textipa{[\u*{i}]} и \textipa{[\u*{u}]}, которые,
объединяясь с рядом стоящими гласными, образуют ложные дифтонги.

Сочетания <<полугласный + гласный>>:

\begin{longtblr}{
    colspec = { X[1,c,t] X[1,c,t] X[1,c,t] X[4,l,t] },
    rowhead=1,
    row{1} = {font=\bfseries},
}
\toprule
Румынский звук & Русский аналог & Написание & Примеры \\
\midrule
\textipa{[\u*{i}a]}
& я
& ia
& iarba \textipa{["\u*{i}arba]}, poiana \textipa{[po"\u*{i}ana]}, oaia \textipa{["\t*{oa}\u*{i}a]}
\\
\textipa{[\u*{i}e]}
& е
& ie, e-
& iepure \textipa{["\u*{i}epure]}, este \textipa{["\u*{i}este]}
\\
\textipa{[\u*{i}o]}
& ё
& io
& iobag \textipa{[\u*{i}o"bag]}, iod \textipa{[\u*{i}od]}
\\
\textipa{[\u*{i}u]}
& ю
& iu
& iute \textipa{["\u*{i}ute]}
\\
\textipa{[\u*{u}a]}
& уа
& ua
& basmaua \textipa{[bas"ma\u*{u}a]}, măseaua \textipa{[m@"s\t*{ea}\u*{u}a]}
\\
\textipa{[\u*{u}@]}
& --
& uă
& rouă \textipa{["ro\u*{u}@]}
\\
\bottomrule
\end{longtblr}

Сочетания <<гласный + полугласный>>:

\begin{longtblr}{
    colspec = { X[1,c,t] X[1,c,t] X[1,c,t] X[1,l,t] },
    rowhead=1,
    row{1} = {font=\bfseries},
}
\toprule
Румынский звук & Русский аналог & Написание & Пример \\
\midrule
\textipa{[a\u*{i}]}
& ай
& ai
& mai \textipa{[ma\u*{i}]}
\\
\textipa{[@\u*{i}]}
& --
& ăi
& măi \textipa{[m@\u*{i}]}
\\
\textipa{[e\u*{i}]}
& ей
& ei
& mei \textipa{[me\u*{i}]}
\\
\textipa{[I\u*{i}]}
& ий
& ii
& copii \textipa{[ko"pI\u*{i}]}
\\
\textipa{[1\u*{i}]}
& ый
& âi
& câine \textipa{["k1\u*{i}ne]}
\\
\textipa{[o\u*{i}]}
& ой
& oi
& apoi \textipa{[a"po\u*{i}]}
\\
\textipa{[u\u*{i}]}
& уй
& ui
& pui \textipa{[pu\u*{i}]}
\\
\textipa{[a\u*{u}]}
& ау
& au
& stau \textipa{[sta\u*{u}]}
\\
\textipa{[@\u*{u}]}
& --
& ău
& tău \textipa{[t@\u*{u}]}
\\
\textipa{[e\u*{u}]}
& эу
& eu
& meu \textipa{[me\u*{u}]}
\\
\textipa{[I\u*{u}]}
& иу
& iu
& viu \textipa{[vI\u*{u}]}
\\
\textipa{[1\u*{u}]}
& ыу
& âu
& brâu \textipa{[br1\u*{u}]}
\\
\textipa{[o\u*{u}]}
& оу
& ou
& bou \textipa{[bo\u*{u}]}
\\
\bottomrule
\end{longtblr}

Кроме того, существуют ложные трифтонги (сочетания дифтонга с полугласным
или гласного с двумя полугласными):

\begin{longtblr}{
    colspec = { X[1,c,t] X[1,c,t] X[1,l,t] },
    rowhead=1,
    row{1} = {font=\bfseries},
}
\toprule
Румынский звук & Написание & Пример \\
\midrule
\textipa{[\t*{ea}\u*{i}]}
& eai
& puneai \textipa{[pu"n\t*{ea}\u*{i}]}
\\
\textipa{[\u*{i}a\u*{i}]}
& iai
& tăiai \textipa{[t@"\u*{i}a\u*{i}]}
\\
\textipa{[\u*{i}e\u*{i}]}
& iei
& miei \textipa{[m\u*{i}e\u*{i}]}
\\
\textipa{[\t*{ea}\u*{u}]}
& eau
& beau \textipa{[b\t*{ea}\u*{u}]}
\\
\textipa{[\u*{i}a\u*{u}]}
& iau
& vroiau \textipa{[vro"\u*{i}a\u*{u}]}
\\
\textipa{[\u*{i}\t*{oa}]}
& ioa
& inimioară \textipa{[InIm"\u*{i}\t*{oa}r@]}
\\
\textipa{[\t*{oa}\u*{i}]}
& oai
& ursoaică \textipa{[ur"s\t*{oa}\u*{i}k@]}
\\
\bottomrule
\end{longtblr}

Местоимение eu я произносится \textipa{[\u*{i}e\u*{u}]},
местоимение ei они --- \textipa{[\u*{i}e\u*{i}]}.
